% GENERATED FILE. Edit canonical lessons, then run npm run generate:books.
% canonical-source-hash: fnv1a64:003597cb9b038427
% canonical-lessons: ML-C500-velampal, ML-C500-maramkotti, ML-C500-desatanam, ML-C500-cekkeruka, ML-C500-kutukuttuka, ML-R500-review-493-500

\chapter{Birds}
\label{ch:ml-birds}
\hlchaptermodality{500}

\begin{chapteropening}
\textbf{By the end of this chapter:} \emph{I can say a hornbill, a woodpecker, migration, to roost and to build a nest.}

Complete the last lesson of chapter 500: I can say a hornbill, a woodpecker, migration, to roost and to build a nest.
\end{chapteropening}

\section[\texorpdfstring{\ml{വേഴാമ്പൽ} (vēḻāmpal)}{വേഴാമ്പൽ (vēḻāmpal)}]{\ml{വേഴാമ്പൽ} (vēḻāmpal) — a hornbill}

\label{lesson:ML-C500-velampal}



\begin{quote}
Before the new one: say the Malayalam for a wild elephant, then the Malayalam for a gaur.
\end{quote}

\subsection*{You'll want to know: \ml{വേഴാമ്പൽ}}
\textbf{\ml{വേഴാമ്പൽ}} — \emph{vēḻāmpal} — \textquotedblleft{}a hornbill\textquotedblright{}.

\begin{quote}
\ml{വേഴാമ്പൽ} \ml{ഉയരമുള്ള} \ml{മരങ്ങളിൽ} \ml{താമസിക്കുന്നു}.
\end{quote}

\begin{quote}
\emph{vēḻāmpal uyaramuḷḷa maraṅṅaḷil tāmasikkunnu.}
\end{quote}

\begin{quote}
The hornbill lives in tall trees.
\end{quote}

A big forest bird with a heavy beak; it is the state bird.

\begin{cousinweb}
\textbf{\ml{വേഴാമ്പൽ}} (\emph{vēḻāmpal}) is a Malayalam bird name whose origin is not settled.
\end{cousinweb}

\begin{grammarlens}[title={What you've built}]
One more word for getting around.
\end{grammarlens}

\subsection*{Your turn}
\begin{itemize}
\raggedright
  \item \emph{Say it:} \emph{vēḻāmpal}
  \item \emph{Say it:} \emph{vēḻāmpal}, once more
  \item \emph{Say it:} \emph{vēḻāmpal uyaramuḷḷa maraṅṅaḷil tāmasikkunnu.} — The hornbill lives in tall trees.
  \item \emph{Recall:} say the Malayalam for a wild elephant, then the Malayalam for a gaur, then say \emph{vēḻāmpal} again
\end{itemize}

\begin{tcolorbox}[breakable,colback=teal!4,colframe=teal!35!black,title={Before you move on}]
What does \ml{വേഴാമ്പൽ} mean? (\textquotedblleft{}A hornbill\textquotedblright{}.) Say it once more.
\end{tcolorbox}
\section[\texorpdfstring{\ml{മരംകൊത്തി} (maraṁkotti)}{മരംകൊത്തി (maraṁkotti)}]{\ml{മരംകൊത്തി} (maraṁkotti) — a woodpecker}

\label{lesson:ML-C500-maramkotti}



\begin{quote}
Before the new one: say the Malayalam for a gaur, then the Malayalam for a hornbill.
\end{quote}

\subsection*{You'll want to know: \ml{മരംകൊത്തി}}
\textbf{\ml{മരംകൊത്തി}} — \emph{maraṁkotti} — \textquotedblleft{}a woodpecker\textquotedblright{}.

\begin{quote}
\ml{ഒരു} \ml{മരംകൊത്തി} \ml{പഴയ} \ml{മരത്തിൽ} \ml{കൊത്തി}.
\end{quote}

\begin{quote}
\emph{oru maraṁkotti paḻaya marattil kotti.}
\end{quote}

\begin{quote}
A woodpecker pecked at the old tree.
\end{quote}

Literally the tree-pecker: \ml{മരം} plus \ml{കൊത്തുക}, to peck, with -\ml{ഇ} for the one who does it.

\begin{cousinweb}
\textbf{\ml{മരംകൊത്തി}} (\emph{maraṁkotti}) is a native Dravidian compound of \ml{മരം}, a tree (DEDR 4711), and \ml{കൊത്തുക}, to peck (DEDR 2091).
\end{cousinweb}

\begin{grammarlens}[title={What you've built}]
One more word for getting around.
\end{grammarlens}

\subsection*{Your turn}
\begin{itemize}
\raggedright
  \item \emph{Say it:} \emph{maraṁkotti}
  \item \emph{Say it:} \emph{maraṁkotti}, once more
  \item \emph{Say it:} \emph{oru maraṁkotti paḻaya marattil kotti.} — A woodpecker pecked at the old tree.
  \item \emph{Recall:} say the Malayalam for a gaur, then the Malayalam for a hornbill, then say \emph{maraṁkotti} again
\end{itemize}

\begin{tcolorbox}[breakable,colback=teal!4,colframe=teal!35!black,title={Before you move on}]
What does \ml{മരംകൊത്തി} mean? (\textquotedblleft{}A woodpecker\textquotedblright{}.) Say it once more.
\end{tcolorbox}
\section[\texorpdfstring{\ml{ദേശാടനം} (dēśāṭanaṁ)}{ദേശാടനം (dēśāṭanaṁ)}]{\ml{ദേശാടനം} (dēśāṭanaṁ) — migration (of birds)}

\label{lesson:ML-C500-desatanam}



\begin{quote}
Before the new one: say the Malayalam for a hornbill, then the Malayalam for a woodpecker.
\end{quote}

\subsection*{You'll want to know: \ml{ദേശാടനം}}
\textbf{\ml{ദേശാടനം}} — \emph{dēśāṭanaṁ} — \textquotedblleft{}migration (of birds)\textquotedblright{}.

\begin{quote}
\ml{തണുപ്പുകാലത്ത്} \ml{പല} \ml{പക്ഷികളും} \ml{ദേശാടനം} \ml{നടത്തുന്നു}.
\end{quote}

\begin{quote}
\emph{taṇuppukālattŭ pala pakṣikaḷuṁ dēśāṭanaṁ naṭattunnu.}
\end{quote}

\begin{quote}
In the cold season many birds migrate.
\end{quote}

A migrating bird is a \ml{ദേശാടനപ്പക്ഷി}. The verb phrase is \ml{ദേശാടനം} \ml{നടത്തുക}.

\begin{cousinweb}
\textbf{\ml{ദേശാടനം}} (\emph{dēśāṭanaṁ}) is a Malayalam compound built from two Sanskrit words, \emph{deśa}, \textquotedblleft{}a country\textquotedblright{}, and \emph{aṭana}, \textquotedblleft{}roaming about\textquotedblright{} (Monier-Williams).
\end{cousinweb}

\begin{grammarlens}[title={What you've built}]
One more word for getting around.
\end{grammarlens}

\subsection*{Your turn}
\begin{itemize}
\raggedright
  \item \emph{Say it:} \emph{dēśāṭanaṁ}
  \item \emph{Say it:} \emph{dēśāṭanaṁ}, once more
  \item \emph{Say it:} \emph{taṇuppukālattŭ pala pakṣikaḷuṁ dēśāṭanaṁ naṭattunnu.} — In the cold season many birds migrate.
  \item \emph{Recall:} say the Malayalam for a hornbill, then the Malayalam for a woodpecker, then say \emph{dēśāṭanaṁ} again
\end{itemize}

\begin{tcolorbox}[breakable,colback=teal!4,colframe=teal!35!black,title={Before you move on}]
What does \ml{ദേശാടനം} mean? (\textquotedblleft{}Migration (of birds)\textquotedblright{}.) Say it once more.
\end{tcolorbox}
\section[\texorpdfstring{\ml{ചേക്കേറുക} (cēkkēṟuka)}{ചേക്കേറുക (cēkkēṟuka)}]{\ml{ചേക്കേറുക} (cēkkēṟuka) — to roost; to come home to roost}

\label{lesson:ML-C500-cekkeruka}



\begin{quote}
Before the new one: say the Malayalam for a woodpecker, then the Malayalam for migration.
\end{quote}

\subsection*{You'll want to know: \ml{ചേക്കേറുക}}
\textbf{\ml{ചേക്കേറുക}} — \emph{cēkkēṟuka} — \textquotedblleft{}to roost; to come home to roost\textquotedblright{}.

\begin{quote}
\ml{സന്ധ്യയ്ക്ക്} \ml{കാക്കകൾ} \ml{മരങ്ങളിൽ} \ml{ചേക്കേറി}.
\end{quote}

\begin{quote}
\emph{sandhyaykkŭ kākkakaḷ maraṅṅaḷil cēkkēṟi.}
\end{quote}

\begin{quote}
At dusk the crows went to roost in the trees.
\end{quote}

Past: \ml{ചേക്കേറി}, \emph{cēkkēṟi}.

\ml{ചേക്ക്}, a roost, plus \ml{ഏറുക}, to climb up. People who settle somewhere new are said to \ml{ചേക്കേറുക} too.

\begin{cousinweb}
\textbf{\ml{ചേക്കേറുക}} (\emph{cēkkēṟuka}) is a native Dravidian compound of \ml{ചേക്ക്}, a roost (DEDR 1990), and \ml{ഏറുക}, to climb up.
\end{cousinweb}

\begin{grammarlens}[title={What you've built}]
One more word for getting around.
\end{grammarlens}

\subsection*{Your turn}
\begin{itemize}
\raggedright
  \item \emph{Say it:} \emph{cēkkēṟuka}
  \item \emph{Say it:} \emph{cēkkēṟuka}, once more
  \item \emph{Say it:} \emph{sandhyaykkŭ kākkakaḷ maraṅṅaḷil cēkkēṟi.} — At dusk the crows went to roost in the trees.
  \item \emph{Recall:} say the Malayalam for a woodpecker, then the Malayalam for migration, then say \emph{cēkkēṟuka} again
\end{itemize}

\begin{tcolorbox}[breakable,colback=teal!4,colframe=teal!35!black,title={Before you move on}]
What does \ml{ചേക്കേറുക} mean? (\textquotedblleft{}To roost; to come home to roost\textquotedblright{}.) Say it once more.
\end{tcolorbox}
\section[\texorpdfstring{\ml{കൂടുകൂട്ടുക} (kūṭukūṭṭuka)}{കൂടുകൂട്ടുക (kūṭukūṭṭuka)}]{\ml{കൂടുകൂട്ടുക} (kūṭukūṭṭuka) — to build a nest}

\label{lesson:ML-C500-kutukuttuka}



\begin{quote}
Before the new one: say the Malayalam for migration, then the Malayalam for to roost.
\end{quote}

\subsection*{You'll want to know: \ml{കൂടുകൂട്ടുക}}
\textbf{\ml{കൂടുകൂട്ടുക}} — \emph{kūṭukūṭṭuka} — \textquotedblleft{}to build a nest\textquotedblright{}.

\begin{quote}
\ml{കുരുവികൾ} \ml{വരാന്തയിൽ} \ml{കൂടുകൂട്ടി}.
\end{quote}

\begin{quote}
\emph{kuruvikaḷ varāntayil kūṭukūṭṭi.}
\end{quote}

\begin{quote}
Sparrows built a nest on the verandah.
\end{quote}

Past: \ml{കൂടുകൂട്ടി}, \emph{kūṭukūṭṭi}.

\ml{കൂട്}, a nest, plus \ml{കൂട്ടുക}, to put together.

\begin{cousinweb}
\textbf{\ml{കൂടുകൂട്ടുക}} (\emph{kūṭukūṭṭuka}) is a native Dravidian compound of \ml{കൂട്}, a nest (DEDR 1883), and \ml{കൂട്ടുക}, to put together (DEDR 1882).
\end{cousinweb}

\begin{grammarlens}[title={What you've built}]
One more word for getting around.
\end{grammarlens}

\subsection*{Your turn}
\begin{itemize}
\raggedright
  \item \emph{Say it:} \emph{kūṭukūṭṭuka}
  \item \emph{Say it:} \emph{kūṭukūṭṭuka}, once more
  \item \emph{Say it:} \emph{kuruvikaḷ varāntayil kūṭukūṭṭi.} — Sparrows built a nest on the verandah.
  \item \emph{Recall:} say the Malayalam for migration, then the Malayalam for to roost, then say \emph{kūṭukūṭṭuka} again
\end{itemize}

\begin{tcolorbox}[breakable,colback=teal!4,colframe=teal!35!black,title={Before you move on}]
What does \ml{കൂടുകൂട്ടുക} mean? (\textquotedblleft{}To build a nest\textquotedblright{}.) Say it once more.
\end{tcolorbox}
\section[(dialogue)]{Words from chapters 493–500 — a review}

\label{lesson:ML-R500-review-493-500}



\begin{quote}
Say the Malayalam for a consumer and a customer.
\end{quote}

\subsection*{Your turn}
Say these aloud:

\begin{itemize}
\raggedright
  \item a consumer, a customer, a warranty, to test, to swap for another
  \item to take back, an expense, to tally, to count and check, to sell off
  \item a peak, a hillside, a mountain pass, a meadow, shade
  \item a large river, a brook, a water spring, a dam, to dry up
  \item the coastal region, an estuary, high tide, low tide, a sea surge
  \item a mangrove forest, a swamp, a coral reef, a shark, to accumulate
  \item a forest, a wild animal, a sanctuary, a wild elephant, a gaur
  \item a hornbill, a woodpecker, migration, to roost, to build a nest
\end{itemize}

\begin{tcolorbox}[breakable,colback=teal!4,colframe=teal!35!black,title={Before you move on}]
A consumer? (\textbf{\ml{ഉപഭോക്താവ്}}, \emph{upabhōktāvŭ}.) A customer? (\textbf{\ml{ഇടപാടുകാരൻ}}, \emph{iṭapāṭukāran}.) A warranty? (\textbf{\ml{വാറന്റി}}, \emph{vāṟanṟi}.) To test? (\textbf{\ml{പരീക്ഷിക്കുക}}, \emph{parīkṣikkuka}.) To swap for another? (\textbf{\ml{മാറ്റിയെടുക്കുക}}, \emph{māṟṟiyeṭukkuka}.) To take back? (\textbf{\ml{തിരിച്ചെടുക്കുക}}, \emph{tiricceṭukkuka}.) An expense? (\textbf{\ml{ചെലവ്}}, \emph{celavŭ}.) To tally? (\textbf{\ml{ഒത്തുനോക്കുക}}, \emph{ottunōkkuka}.) To count and check? (\textbf{\ml{എണ്ണിനോക്കുക}}, \emph{eṇṇinōkkuka}.) To sell off? (\textbf{\ml{വിറ്റഴിക്കുക}}, \emph{viṟṟaḻikkuka}.) A peak? (\textbf{\ml{കൊടുമുടി}}, \emph{koṭumuṭi}.) A hillside? (\textbf{\ml{മലഞ്ചെരിവ്}}, \emph{malañcerivŭ}.) A mountain pass? (\textbf{\ml{ചുരം}}, \emph{curaṁ}.) A meadow? (\textbf{\ml{പുൽമേട്}}, \emph{pulmēṭŭ}.) Shade? (\textbf{\ml{തണൽ}}, \emph{taṇal}.) A large river? (\textbf{\ml{നദി}}, \emph{nadi}.) A brook? (\textbf{\ml{അരുവി}}, \emph{aruvi}.) A water spring? (\textbf{\ml{ഉറവ}}, \emph{uṟava}.) A dam? (\textbf{\ml{അണക്കെട്ട്}}, \emph{aṇakkeṭṭŭ}.) To dry up? (\textbf{\ml{വറ്റുക}}, \emph{vaṟṟuka}.) The coastal region? (\textbf{\ml{തീരദേശം}}, \emph{tīradēśaṁ}.) An estuary? (\textbf{\ml{അഴിമുഖം}}, \emph{aḻimukhaṁ}.) High tide? (\textbf{\ml{വേലിയേറ്റം}}, \emph{vēliyēṟṟaṁ}.) Low tide? (\textbf{\ml{വേലിയിറക്കം}}, \emph{vēliyiṟakkaṁ}.) A sea surge? (\textbf{\ml{കടലാക്രമണം}}, \emph{kaṭalākramaṇaṁ}.) A mangrove forest? (\textbf{\ml{കണ്ടൽക്കാട്}}, \emph{kaṇṭalkkāṭŭ}.) A swamp? (\textbf{\ml{ചതുപ്പ്}}, \emph{catuppŭ}.) A coral reef? (\textbf{\ml{പവിഴപ്പുറ്റ്}}, \emph{paviḻappuṟṟŭ}.) A shark? (\textbf{\ml{സ്രാവ്}}, \emph{srāvŭ}.) To accumulate? (\textbf{\ml{അടിഞ്ഞുകൂടുക}}, \emph{aṭiññukūṭuka}.) A forest? (\textbf{\ml{വനം}}, \emph{vanaṁ}.) A wild animal? (\textbf{\ml{വന്യജീവി}}, \emph{vanyajīvi}.) A sanctuary? (\textbf{\ml{സങ്കേതം}}, \emph{saṅkētaṁ}.) A wild elephant? (\textbf{\ml{കാട്ടാന}}, \emph{kāṭṭāna}.) A gaur? (\textbf{\ml{കാട്ടുപോത്ത്}}, \emph{kāṭṭupōttŭ}.) A hornbill? (\textbf{\ml{വേഴാമ്പൽ}}, \emph{vēḻāmpal}.) A woodpecker? (\textbf{\ml{മരംകൊത്തി}}, \emph{maraṁkotti}.) Migration? (\textbf{\ml{ദേശാടനം}}, \emph{dēśāṭanaṁ}.) To roost? (\textbf{\ml{ചേക്കേറുക}}, \emph{cēkkēṟuka}.) To build a nest? (\textbf{\ml{കൂടുകൂട്ടുക}}, \emph{kūṭukūṭṭuka}.)
\end{tcolorbox}
